\section{Oracle Headroom at Fixed Memory (R2-02)}
\label{sec:oracle-headroom}

Having established replay fidelity (Section~\ref{sec:replay-validation}),
we next compare four retention policies at a single, fixed 4~GiB
expert-memory budget -- the capture budget of the validated trace -- over
the same decode-phase workload:
\begin{itemize}
 \item \textbf{Runtime-Sharded-LRU}: a statistically shard-matched LRU
   (16 shards, the real runtime's own shard count at 4~GiB), causal.
   Its shard-to-key mapping is statistical, not the runtime's exact
   address-derived mapping (Limitation, Section~\ref{sec:limitations}).
 \item \textbf{Global-LRU}: an unsharded, global-scope LRU, causal.
 \item \textbf{LFU}: an unsharded, global-scope least-frequently-used
   policy, causal. Treated as a secondary baseline only throughout this
   report.
 \item \textbf{Oracle-Global-Belady}: the non-causal upper bound
   (Section~\ref{sec:problem-formulation}).
\end{itemize}

\newcommand{\fixthis}{}
\begin{table}[t]
\centering
\caption{Policy comparison at the 4~GiB anchor budget (R-02).}
\label{tab:r2-02-comparison}
\small
\begin{tabular}{@{}lrrrr@{}}
\toprule
Policy & Causal & Misses & Miss rate & Miss bytes (GB) \\
\midrule
Runtime-Sharded-LRU & yes & 35,943 & 0.2457 & 103.35 \\
Global-LRU & yes & 35,529 & 0.2428 & 102.15 \\
LFU & yes & 39,124 & 0.2674 & 112.12 \\
Oracle-Global-Belady & no & 14,754 & 0.1008 & 42.30 \\
\bottomrule
\end{tabular}
\end{table}


\paragraph{Primary comparison.} Table~
ef{tab:r2-02-comparison}
reports the full four-policy comparison. The primary, scientifically
clean comparison is Global-LRU vs.\ Oracle-Global-Belady -- both at
global cache scope, differing only in causal information availability.
At 4~GiB, Oracle-Global-Belady reduces miss bytes relative to Global-LRU
by a relative amount of $0.5859$ (an absolute reduction of
$\mathbf{59{,}855{,}155{,}200}$ bytes), i.e.\ an information headroom
$H_{4\text{GiB}}(\text{Global-LRU}) \approx 58.59\%$
(Definition~\ref{def:information-headroom}).

\paragraph{Structural sharding gap.} The Runtime-Sharded-LRU vs.\
Global-LRU comparison isolates the cost of the runtime's own sharded
retention architecture, holding causal information availability fixed.
At 4~GiB this structural gap is small and, notably, of the \emph{opposite
sign} at lower budgets than at this one (Section~\ref{sec:recoverable-envelope}):
at 4~GiB, Runtime-Sharded-LRU incurs $1{,}200{,}218{,}112$ more miss bytes
than Global-LRU (a relative structural cost of $1.16\%$), i.e.\ sharding
itself is mildly disadvantageous relative to a global LRU at this
specific budget. No Research Part~1 shard-constrained oracle exists, so
the structural and policy-information gaps are decomposed using
Global-LRU as a shared pivot (Section~\ref{sec:recoverable-envelope});
this is a stated limitation, not hidden.

\paragraph{LFU.} LFU underperforms Global-LRU at this budget (higher
miss count and miss bytes) and is treated throughout this report as a
secondary causal baseline only; it plays no role in the structural/
information decomposition.
